\chapter{Parámetros de Calidad de la Fuente}
\label{sec.c5:parametros_calidad}

La evaluación integral del desempeño de una fuente regulada se cuantifica mediante los factores de calidad
definidos en el marco teórico de la cátedra \cite{gonzalez_fuentes2}. A continuación se detallan las
definiciones formales, las mediciones experimentales y el cálculo de cada parámetro para las condiciones de
diseño exigidas ($I_{Lmin} = \SI{100}{\milli\ampere}$, $I_{Lmax} = \SI{3}{\ampere}$).

\section{Regulación de carga ($R_{car}$)}
\label{sec.c5:rcar}

La regulación de carga expresa la variación porcentual de la tensión de salida cuando la corriente de carga
pasa de su valor mínimo especificado a plena carga (Ecuación 6-42 de \cite{gonzalez_fuentes2}):
\begin{equation}
    R_{car} = \frac{V_{o(NL)} - V_{o(FL)}}{V_{o(nom)}} \times 100
    \label{eq.c5:rcar}
\end{equation}
A partir de los datos registrados en el archivo \lstinline{vRL_120_4.txt}:
\begin{itemize}
    \item A corriente mínima ($I_{Lmin} = \SI{100}{\milli\ampere}$, $R_L = \SI{120}{\ohm}$):
    $V_{o(NL)} = \SI{12.00156}{\volt}$.
    \item A plena carga nominal ($I_{Lmax} = \SI{3.0}{\ampere}$, $R_L = \SI{4}{\ohm}$):
    $V_{o(FL)} = \SI{11.99944}{\volt}$.
\end{itemize}
La variación absoluta es de apenas $\Delta V_o = \SI{2.12}{\milli\volt}$. Aplicando la ecuación~\eqref{eq.c5:rcar}:
\begin{equation}
    R_{car} = \frac{\SI{12.00156}{\volt} - \SI{11.99944}{\volt}}{\SI{12.00}{\volt}} \times 100 = \SI{0.0177}{\percent}
    \label{eq.c5:rcar_num}
\end{equation}

\section{Resistencia interna de salida ($R_{sal}$)}
\label{sec.c5:rsal}

La resistencia dinámica de salida modela la pendiente de caída de tensión frente a incrementos en la demanda
de corriente (Ecuación 6-43 de \cite{gonzalez_fuentes2}):
\begin{equation}
    R_{sal} = \frac{\Delta V_o}{\Delta I_L} = \frac{V_{o(NL)} - V_{o(FL)}}{I_{Lmax} - I_{Lmin}}
    \label{eq.c5:rsal}
\end{equation}
Reemplazando los valores experimentales:
\begin{equation}
    R_{sal} = \frac{\SI{2.12}{\milli\volt}}{\SI{3.00}{\ampere} - \SI{0.10}{\ampere}}
            = \frac{\SI{2.12}{\milli\volt}}{\SI{2.90}{\ampere}}
            = \mathbf{\SI{0.731}{\milli\ohm}}
    \label{eq.c5:rsal_num}
\end{equation}
El valor obtenido en el orden de los submiliohmios confirma el comportamiento cuasi-ideal del lazo de regulación
del integrado LM7812 asistido por el booster exterior.

\section{Regulación de red o línea ($R_{red}$)}
\label{sec.c5:rred}

La regulación de red mide la capacidad de la fuente para aislar la salida frente a fluctuaciones en la red eléctrica
de entrada no regulada (Ecuación 6-41 de \cite{gonzalez_fuentes2}):
\begin{equation}
    R_{red} = \frac{\Delta V_o}{\Delta V_{in}} \times 100
    \label{eq.c5:rred}
\end{equation}
Para evaluar una tolerancia de línea del $\pm 10\%$ sobre la tensión nominal de entrada ($V_{in} = \SI{18}{\volt}$):
\begin{itemize}
    \item $V_{in1} = 18 - 10\% = \SI{16.2}{\volt} \implies V_o = \SI{11.99944}{\volt}$
    (datos de \lstinline{vRl_vin_variable_regulacion_red.txt}).
    \item $V_{in2} = 18 + 10\% = \SI{19.8}{\volt} \implies V_o = \SI{11.99945}{\volt}$.
\end{itemize}
Ante una variación de entrada de $\Delta V_{in} = \SI{3.6}{\volt}$, la tensión de salida experimenta un corrimiento
de solo $\Delta V_o = \SI{10}{\micro\volt}$. Por ende:
\begin{equation}
    R_{red} = \frac{\SI{10}{\micro\volt}}{\SI{3.6}{\volt}} \times 100 = \mathbf{\SI{0.00028}{\percent}}
    \quad \left(\SI{0.0028}{\milli\volt/\volt}\right)
    \label{eq.c5:rred_num}
\end{equation}

\section{Coeficiente térmico de tensión ($C.T.$)}
\label{sec.c5:ct}

El coeficiente de temperatura cuantifica el desplazamiento térmico de la referencia interna bandgap y del lazo
amplificador (Ecuación 6-44 de \cite{gonzalez_fuentes2}):
\begin{equation}
    C.T. = \frac{\Delta V_o}{\Delta T}
    \label{eq.c5:ct}
\end{equation}
A partir del ensayo de temperatura (\lstinline{vRL_temp_0_25_55.txt}) a plena carga nominal ($R_L = \SI{4}{\ohm}$):
\begin{itemize}
    \item A $T_1 = \SI{0}{\celsius}$: $V_o(\SI{0}{\celsius}) = \SI{12.04379}{\volt}$.
    \item A $T_2 = \SI{25}{\celsius}$: $V_o(\SI{25}{\celsius}) = \SI{12.00316}{\volt}$.
\end{itemize}
Para el intervalo $\Delta T = \SI{25}{\celsius}$:
\begin{equation}
    C.T. = \frac{\SI{12.00316}{\volt} - \SI{12.04379}{\volt}}{\SI{25}{\celsius}}
         = \frac{\SI{-40.63}{\milli\volt}}{\SI{25}{\celsius}}
         = \mathbf{\SI{-1.625}{\milli\volt/\celsius}}
    \label{eq.c5:ct_num}
\end{equation}
Normalizado a la tensión nominal de $\SI{12}{\volt}$, representa un coeficiente relativo de:
\begin{equation*}
    C.T._{rel} = \frac{\SI{-1.625}{\milli\volt/\celsius}}{\SI{12}{\volt}} \times 10^6 \approx \SI{-135.4}{ppm/\celsius}
\end{equation*}

\section{Tensión de ripple y Rechazo de rizado ($RR$)}
\label{sec.c5:ripple}

Para determinar la atenuación del zumbido de rectificación, se inyectó a la entrada una señal senoidal de
$100\text{ Hz}$ con amplitud de $\SI{1}{\volt}$ (equivalente a $V_{rip(in)} = \SI{2.0}{\volt_{pp}}$) montada sobre el nivel
continuo de $\SI{18}{\volt}$ (\lstinline{V1 N001 0 SINE(18 1 100)}), operando a plena carga ($I_L = \SI{3}{\ampere}$).

\begin{figure}[!ht]
    \centering
    \begin{tikzpicture}
\begin{axis}[
    width=0.85\textwidth,
    height=5.8cm,
    xlabel={Tiempo en régimen permanente $t$ [\si{\milli\second}]},
    ylabel={Rizado $\Delta V_{out}(t)$ [\si{\micro\volt}]},
    grid=both,
    grid style={line width=.1pt, draw=gray!20},
    major grid style={line width=.2pt, draw=gray!40},
    xmin=0, xmax=10,
    ymin=-6, ymax=6,
    legend pos=north east,
    legend style={font=\small}
]

% Oscilacion de 100 Hz centrada
\addplot[
    color=teal!80!black,
    line width=1.3pt,
    smooth
] coordinates {
    (0.00, 0.00)
    (0.34, 1.47)
    (1.12, 3.38)
    (2.19, 4.33)
    (3.32, 2.42)
    (4.19, 0.52)
    (5.00, -1.20)
    (6.10, -3.30)
    (6.93, -4.25)
    (7.42, -4.25)
    (8.25, -3.30)
    (9.22, -1.39)
    (10.00, 0.00)
};
\addlegendentry{Rizado a plena carga ($I_L = \SI{3}{\ampere}$)}

\draw[<->, >=stealth, line width=0.8pt, color=red!70!black] 
    (axis cs: 2.19, 4.33) -- (axis cs: 2.19, -4.25)
    node[midway, right, font=\footnotesize] {$V_{rpp} = \SI{8.58}{\micro\volt}$};

\end{axis}
\end{tikzpicture}

    \caption{Ondulación residual de tensión de salida (ripple de $\SI{100}{\hertz}$) a plena carga nominal.}
    \label{fig.c5:ripple}
\end{figure}

Del registro transitorio en régimen estacionario (\lstinline{vRl_vin_variable_ripple.txt}), ilustrado en la
figura~\ref{fig.c5:ripple}, se extrae una tensión de rizado pico a pico residual en la carga de:
\begin{equation}
    V_{rpp(out)} = \SI{8.58}{\micro\volt}
    \label{eq.c5:vrpp}
\end{equation}
De acuerdo a la Ecuación 6-78 de la cátedra \cite{gonzalez_fuentes2}, la tensión eficaz de ripple para una
señal de rizado se calcula como:
\begin{equation}
    V_{ef-rip} = \frac{V_{rpp(out)}}{3.464} = \frac{\SI{8.58}{\micro\volt}}{3.464} = \mathbf{\SI{2.48}{\micro\volt_{rms}}}
    \label{eq.c5:vef_rip}
\end{equation}
El factor de rechazo de ripple ($RR$) de la fuente resulta:
\begin{equation}
    RR = 20 \log_{10} \left( \frac{V_{rip\_in(pp)}}{V_{rpp(out)}} \right)
       = 20 \log_{10} \left( \frac{\SI{2.0}{\volt}}{\SI{8.58}{\micro\volt}} \right)
       = \mathbf{\SI{107.35}{\decibel}}
    \label{eq.c5:rr}
\end{equation}
Este nivel de rechazo supera holgadamente los $\SI{60}{\decibel}$ especificados típicamente en la hoja de datos
del LM7812 \cite{ti_lm78xx}, producto del efecto combinado del regulador integrado y los capacitores de desacoplo.

\section{Tabla consolidada de factores de calidad}
\label{sec.c5:tabla_calidad}

En la tabla~\ref{tab.c5:calidad} se compendian todos los parámetros de calidad obtenidos para la fuente diseñada.

\begin{table}[!ht]
\centering
\caption{Resumen de factores de calidad obtenidos de la fuente regulada.}
\label{tab.c5:calidad}
\begin{tabular}{lccc}
\hline
Parámetro de Calidad & Símbolo & Valor Obtenido & Unidad \\
\hline
Tensión nominal de salida & $V_{out}$ & $12.00$ & \si{\volt} \\
Corriente máxima de diseño & $I_{Lmax}$ & $3.00$ & \si{\ampere} \\
Regulación de carga & $R_{car}$ & $0.0177$ & \si{\percent} \\
Resistencia dinámica de salida & $R_{sal}$ & $0.731$ & \si{\milli\ohm} \\
Regulación de red / línea & $R_{red}$ & $0.00028$ & \si{\percent} \\
Factor de regulación de línea & $F_{red}$ & $0.0028$ & \si{\milli\volt/\volt} \\
Coeficiente de temperatura & $C.T.$ & $-1.625$ & \si{\milli\volt/\celsius} \\
Tensión de ripple pico a pico & $V_{rpp}$ & $8.58$ & \si{\micro\volt} \\
Tensión de ripple eficaz & $V_{ef-rip}$ & $2.48$ & \si{\micro\volt_{rms}} \\
Rechazo de ripple ($100\text{ Hz}$) & $RR$ & $107.35$ & \si{\decibel} \\
\hline
\end{tabular}
\end{table}
